\section{Discussion}\label{sec:discussion}

\subsection{When is a daily azimuth estimate trustworthy?}
\draftnote{Summarize practical acceptance criteria, uncertainty calibration, and minimum clear-day requirements. Separate recommendations supported by controlled data from extrapolations.}

\subsection{Identifiability and potential confounding factors}
\draftnote{Address timestamp shifts, incorrect tilt, horizon obstructions, partial shading, directional/cosine response, soiling, atmospheric variability, low curvature, model dependence, and seasonal effects. Clarify the meaning of `effective azimuth'.}

\subsection{Relation to earlier work and scope of the contribution}
The present work builds on the previously published peak-hour-angle method
\citep{piccioli2021} and its field application in a thesis
\citep{piccioli2022thesis}. Its intended contribution is a more rigorous
experimental validation and an explicit reliability assessment, rather than
re-introduction of the original geometric method.
\draftnote{Compare fairly with Barbier et al. and Meng et al. State differences in required inputs, data type, scope, and dependence on transposition models. Make clear what results are genuinely new.}

\subsection{Implications for solar-resource assessment}
\draftnote{Explain when orientation diagnosis can improve metadata QC and model assessment. Avoid conflating reduced model error with proof that the physical azimuth is exact.}
